% =====================================================================
%  Conclusion générale
% =====================================================================
\chapter*{Conclusion générale}
\addcontentsline{toc}{chapter}{Conclusion générale}
\markboth{Conclusion générale}{}

Ce rapport a montré comment l'intelligence artificielle peut quitter les
centres de données pour s'installer au cœur des objets connectés. Partant de
l'architecture classique de l'IoT, où l'objet se contente de mesurer et de
transmettre, nous avons vu que l'AIoT donne aux systèmes connectés la capacité
d'interpréter leurs données, et que le TinyML pousse cette logique jusqu'aux
microcontrôleurs les plus modestes.

Cette évolution repose sur une séparation nette entre l'entraînement, réalisé
sur des machines puissantes, et l'inférence, seule exécutée dans l'objet. Elle
n'est possible que grâce à un pipeline de développement rigoureux et à des
techniques de compression (quantification, élagage, distillation) capables de
réduire la taille d'un modèle d'un ou deux ordres de grandeur, ainsi qu'à des
microcontrôleurs et des outils logiciels de plus en plus adaptés.

Les bénéfices sont réels : réactivité accrue, fonctionnement en l'absence de
réseau, réduction des données transmises et donc meilleure maîtrise de la
confidentialité et, souvent, de la consommation énergétique. Ils doivent
cependant être nuancés : la latence n'est pas nulle, le calcul local consomme
lui aussi de l'énergie, la fiabilité dépend fortement des données
d'entraînement et l'évolution des modèles déployés reste délicate.

Le TinyML ne remplace donc pas le cloud : il le complète. L'avenir de l'AIoT se
dessine comme une intelligence distribuée, dans laquelle chaque niveau (objet,
passerelle, cloud) exécute les traitements qui lui conviennent le mieux. Notre
exemple de détection d'anomalies vibratoires sur ESP32 illustre bien cette
répartition : l'objet surveille et alerte en temps réel, tandis que le cloud
supervise, archive et améliore le modèle au fil du temps.
